\documentclass[12pt]{article}
\usepackage[utf8]{inputenc}
\usepackage[T1]{fontenc}
\usepackage{hyperref}
\title{Tunebo, Western Language Resources for K-12 Curriculum Development}
\author{Compiled by Claude Sonnet 5 (claude-sonnet-5) for Andr\'as Kornai}
\date{}
\begin{document}
\maketitle

\section{Introduction}
Western Tunebo (ISO 639-3 tnb) is a Chibchan language of Colombia. Search results associate it with the Aguas Blancas (Rikuwa) variety. It is closely related to Central Tunebo (tuf), and the two are often grouped as Uw Cuwa (U'wa). Most of the resources found describe Central Tunebo, so they may not match Western Tunebo. This report rests on a single round of web searches.

\section{Primary Educational Resources}
The Colombian Ministry of Culture document on the U'wa mentions teaching materials: a bilingual Tunebo--Spanish dictionary with a brief grammar, and an 11-lesson introductory course. These are not specific to the Western variety: \url{https://www.mincultura.gov.co:443/areas/poblaciones/noticias/Documents/U%C2%B4wa.pdf}

\section{Linguistic and Reference Materials}
\begin{itemize}
\item Headland (1997), bilingual Uw Cuwa--Spanish dictionary with a grammar, hosted on Webonary. It is based mainly on Central Tunebo (Cobar\'ia and Tegr\'ia): \url{https://www.webonary.org/tunebo/overview/introduction/?lang=en}
\item Rochereau (1926, 1927), \emph{La lengua tuneba y sus dialectos}, cited in Glottolog: \url{https://glottolog.org/resource/reference/id/12914}
\item Overview: \url{https://en.wikipedia.org/wiki/Uwa_language}
\end{itemize}

\section{Religious and Community Resources}
None found.

\section{Digital and Audio Resources}
The Webonary dictionary is the only digital resource found. No audio or video was found.

\section{Curriculum Development Context}
Existing materials centre on the Central variety. Any curriculum for Western Tunebo would need to confirm with the community how far the Central materials can be reused. No school programme specific to Western Tunebo was found.
\end{document}